\section*{Chapitre \ref{ch:nw:on-chain-connectivity} --- Connectivité on-chain}

\begin{solution}{exo:nw:on-chain-connectivity:1}
$4 \times 300 = 1\,200$~ms, soit 1.2~seconde~; mais le temps de slot effectif change dynamiquement, il faut donc le lire sur la chaîne plutôt que de supposer la valeur par défaut.
\end{solution}

\begin{solution}{exo:nw:on-chain-connectivity:2}
Un réseau de propagation achemine la transaction plus vite et avec moins de variance que le bavardage public, et peut encore la livrer au public~; un relais privé
la garde hors du mempool public jusqu'à ce qu'un constructeur de blocs l'inclue. Le premier vend de la latence, le second de l'invisibilité.
\end{solution}

\begin{solution}{exo:nw:on-chain-connectivity:3}
Ses pairs sont aléatoires~: un nœud situé dans la même ville est en général atteint par l'intermédiaire d'autres nœuds, chacun ajoutant son traitement, et dans le modèle le nœud médian à Tokyo est
atteint après huit sauts, \qty{40}{\milli\second} contre \qty{1}{\milli\second} via le moteur de Tokyo.
\end{solution}

\begin{solution}{exo:nw:on-chain-connectivity:4}
La diffusion en racine carrée remplace certaines poussées par une annonce suivie d'une demande, ce qui allonge ces liaisons d'un aller-retour~; d'après la proposition, des liaisons plus longues
ne peuvent raccourcir aucun plus court chemin. On l'utilise parce qu'elle envoie chaque transaction en entier à bien moins de pairs, ce qui économise la bande passante au prix d'un léger surcoût en temps.
\end{solution}

\begin{solution}{exo:nw:on-chain-connectivity:5}
Le moteur de Singapour~: \qty{39.9}{\milli\second} contre \qty{58.9}{\milli\second} par bavardage, soit \qty{19}{\milli\second} de moins.
\end{solution}

\begin{solution}{exo:nw:on-chain-connectivity:6}
Un dixième de seconde de propagation représente moins d'un pour cent d'un slot de 12~secondes~: presque toute transaction envoyée tôt dans le slot arrive à temps, où que se trouve
le proposeur, et la géographie ne compte que pour ce qui est envoyé en fin de slot.
\end{solution}

\begin{solution}{exo:nw:on-chain-connectivity:7}
La moitié des nœuds après \qty{80.3}{\milli\second} et neuf dixièmes après \qty{89.5}{\milli\second}~: 4.7 et \qty{10.0}{\milli\second} plus rapidement, car davantage
de pairs signifient moins de sauts.
\end{solution}

\begin{solution}{exo:nw:on-chain-connectivity:8}
Le chemin le plus rapide mène au prochain leader, et non à la société~: un moteur situé près des serveurs de la société n'aide que si le leader est lui aussi proche. Avec des leaders
répartis sur plusieurs continents, l'endpoint doit être choisi pour chaque leader à partir du calendrier.
\end{solution}

\begin{solution}{pb:nw:on-chain-connectivity:1}
\textbf{Partie I.}
\begin{enumerate}
\item \qty{85.0}{\milli\second} et \qty{99.5}{\milli\second}.
\item \qty{113.9}{\milli\second} et \qty{149.0}{\milli\second}.
\item Après \qty{83.5}{\milli\second}.
\item Il faut trois sauts, chacun avec \qty{5}{\milli\second} de traitement, alors que le chemin direct ne compte qu'un moteur à \qty{1}{\milli\second}.
\end{enumerate}
\textbf{Partie II.}
\begin{enumerate}[resume]
\item Celui de Francfort~; \qty{69.5}{\milli\second}.
\item \qty{45.6}{\milli\second} dans un sens~; le modèle multiplie ce temps par~1.5.
\item \qty{1}{\milli\second}, une hypothèse~; le reste provient du tracé.
\item La transaction elle-même~: le bavardage la montre à chaque nœud sur le trajet, le moteur seulement à lui-même et au leader.
\end{enumerate}
\textbf{Partie III.}
\begin{enumerate}[resume]
\item 61~\% avec le moteur, 48~\% par bavardage.
\item 80.5~\% contre 74.0~\%.
\item Dès que le temps restant dépasse l'arrivée la plus lente (\qty{99.5}{\milli\second} par bavardage), les deux méthodes atteignent chaque leader~; l'écart fixe en
  millisecondes représente alors une part plus petite de la fenêtre.
\item Le même leader tient quatre slots, si bien que l'endpoint ne change qu'aux passages de relais~; les transactions envoyées vers la fin de la période d'un leader vont vers la région du leader suivant,
  choisie à l'avance à partir du calendrier.
\end{enumerate}
\textbf{Partie IV.}
\begin{enumerate}[resume]
\item \emph{Résultat nommé}~: dans le modèle, une transaction émise de Tokyo atteint la moitié du réseau en \qty{85.0}{\milli\second} et neuf dixièmes en
  \qty{99.5}{\milli\second} par inondation (113.9 et 149.0 avec la diffusion en racine carrée)~; envoyer au \omterm{def:nw:on-chain-connectivity:engine}{moteur de blocs} le plus proche du prochain leader, plutôt qu'à un nœud public
  à Tokyo, fait passer la part de leaders atteints à temps de 48~\% à 61~\% sur une fenêtre de \qty{150}{\milli\second}, et atteint un leader de Francfort
  en 69.5 au lieu de \qty{83.5}{\milli\second}.
\item Documentation~: les villes et les endpoints des moteurs, les règles de slot et de leader~; hypothèses~: le graphe des pairs, le \omterm{def:nw:the-north-american-map:factor}{facteur de tracé}, les temps de traitement.
\item Des pairs choisis parmi les nœuds proches raccourciraient chaque saut, mais ajouteraient des sauts par-dessus les océans~; les temps du bavardage changeraient, pas le coût de ses sauts supplémentaires.
\item Près des moteurs des régions où se trouvent la plupart des leaders, avec un tracé rapide vers les autres~; pour cette chaîne, l'Europe et un site en Asie.
\item Un second chemin, à variance plus faible, pour ce que le moteur ne transporte pas, et la livraison aux constructeurs de blocs et aux validateurs hors de portée du moteur.
\item Envoyer des transactions marquées par chaque chemin et mesurer le slot dans lequel chacune atterrit, par région de leader~; comparer aux arrivées du modèle.
\item Dans les lignes on-chain~: les moteurs par région, le sélecteur, les relais privés et les réseaux de propagation, et les nœuds que la société fait tourner.
\item Parce que la localisation du prochain leader, connue à l'avance, détermine l'endpoint le plus rapide, et qu'elle change toutes les quelques slots.
\end{enumerate}
\end{solution}

\begin{solution}{iq:nw:on-chain-connectivity:1}
Par bavardage~: chaque nœud le retransmet, ou l'annonce, à ses pairs dès la première réception, après l'avoir validé. C'est lent parce que les pairs sont aléatoires, donc le
chemin zigzague, et chaque saut ajoute validation et file d'attente~; le temps croît avec le nombre de sauts et avec leurs liaisons les plus lentes.
\iqlookfor{Sauts pair à pair~; traitement par saut~; annonce et demande~; queues mesurées.}
\end{solution}

\begin{solution}{iq:nw:on-chain-connectivity:2}
Un service proche des validateurs, qui reçoit les transactions et les bundles et les transmet directement au producteur de blocs, souvent avec une enchère pour l'inclusion.
Il y en a plusieurs parce que les producteurs sont répartis dans le monde et que le dernier saut doit être court où que se trouve le suivant.
\iqlookfor{Transmission directe~; pourboires ou enchère~; endpoints régionaux.}
\end{solution}

\begin{solution}{iq:nw:on-chain-connectivity:3}
On lit le calendrier des prochains slots, on associe chaque leader à une région (publiée ou mesurée), et on envoie chaque transaction vers l'endpoint le plus proche du leader
qui produira encore au moment de son arrivée~; on garde les connexions chaudes dans chaque région, et on mesure les taux d'atterrissage par leader.
\iqlookfor{Calendrier à l'avance~; association des validateurs à des lieux~; passages de relais~; mesure.}
\end{solution}

\begin{solution}{iq:nw:on-chain-connectivity:4}
Quand le nœud se trouve sur le chemin de trading~: un état qui doit être lu avec la latence la plus faible, des transactions que la société ne veut pas partager avec un fournisseur,
ou des taux de requêtes au-delà des forfaits d'un fournisseur~; un endpoint hébergé convient à la lecture de l'historique et à tout ce qui n'est pas une course.
\iqlookfor{Latence~; limites~; confidentialité~; coût d'exploitation.}
\end{solution}

\begin{solution}{iq:nw:on-chain-connectivity:5}
Pour la tenir hors de vue des robots qui passeraient des ordres avant elle ou autour d'elle (frontrunning, sandwiching), et parfois pour obtenir un remboursement de la valeur qu'elle crée.
\iqlookfor{Visibilité~; MEV~; flux d'ordres privé.}
\end{solution}

\begin{solution}{iq:nw:on-chain-connectivity:6}
Des connexions aux \omterm{def:nw:on-chain-connectivity:engine}{moteurs de blocs} ou aux constructeurs de blocs de chaque région de leader, un sélecteur piloté par le \omterm{def:nw:on-chain-connectivity:leader}{calendrier des leaders} avec une marge aux passages de relais, des relais privés
là où la visibilité compte, les nœuds propres de la société dans les grandes régions pour l'état, des nœuds hébergés comme lecteurs de secours, des limiteurs de débit par région, et des taux d'atterrissage
mesurés par région de leader et par chemin.
\iqlookfor{Routage piloté par le calendrier~; présence régionale~; chemins privés~; mesure et limites.}
\end{solution}
